% Hafta 6 — RASP döngüsü: algıla -> savun -> caydır
% Derleme: py -3.12 tools/latex/derle.py docs/week-6/assets/h06-01-rasp-dongusu.tex
\documentclass[tikz,border=8pt]{standalone}
\usepackage{cen429-cizim}
\begin{document}
\begin{tikzpicture}
  \node[baslik] (b) at (0,0) {RASP döngüsü: algıla $\rightarrow$ savun $\rightarrow$ caydır};
  \node[altbaslik] at (b.south west) {koruma uygulamanın İÇİNDE çalışır};
  \node[kutu=38mm, anchor=north west] (n1) at (0,-2.4)
    {\kb{ALGILA}\\bütünlük denetimi\\debugger · kanca\\ortam · root · imza};
  \node[iyikutu=38mm, right=24mm of n1] (n2)
    {\kb{SAVUN}\\fail-closed\\sırrı sil · cihaza bağla\\akış sayacı};
  \node[uyarikutu=38mm, right=24mm of n2] (n3)
    {\kb{CAYDIR}\\decoy çıktı · gecikme\\telemetri\\attestation};
  \draw[ok] (n1.east) -- (n2.west);
  \draw[ok] (n2.east) -- (n3.west);
  \node[fit=(n1)(n2)(n3), inner sep=0pt] (grp) {};
  \node[dip, text width=155mm, anchor=north west] at ($(grp.south west)+(0,-8mm)$)
    {Tek denetim değil, birbirini kontrol eden ÖRTÜŞEN AĞ $\rightarrow$ tek yama yetmez.};
\end{tikzpicture}
\end{document}
